\subsection{円筒による分割}

同じ立体でも,回転軸に垂直に切る代わりに,軸に平行な薄い帯を回す方法がある.たとえば $0\leq a\leq x\leq b$,\,$g(x)\leq y\leq f(x)$ を $y$ 軸のまわりに回す.幅 $\Delta x$ の帯は,半径がほぼ $x$,高さがほぼ $f(x)-g(x)$ の薄い円筒の層になる.

半径 $x$ と $x+\Delta x$ の間の円環の面積は
\[
\pi\{(x+\Delta x)^2-x^2\}
 =2\pi x\Delta x+\pi(\Delta x)^2
\]
である.分割を細かくすると二次の項による誤差の和は零に近づくので,体積は
\begin{equation}\label{式:円筒殻の体積}
V=2\pi\int_a^b x\{f(x)-g(x)\}\dd x
\end{equation}
となる.これを円筒殻による方法という.\,$2\pi x$ は円周の長さ,\,$f-g$ は層の高さである.

\begin{example}{体積の別解}
$0\leq x\leq1$,\,$x^2\leq y\leq x$ の領域を $y$ 軸のまわりに回転させた立体の体積を求めよ.
\end{example}

縦の帯を回すと,半径 $x$,高さ $x-x^2$ の円筒殻になるから,
\[
V=2\pi\int_0^1x(x-x^2)\dd x
 =2\pi\left[\frac{x^3}{3}-\frac{x^4}{4}\right]_0^1
 =\frac{\pi}{6}.
\]
横に切る方法でも確かめられる.\,$0\leq y\leq1$ で,領域の左右の端は $x=y$ と $x=\sqrt y$ である.したがって円環の外半径は $\sqrt y$,内半径は $y$ となり,
\[
V=\pi\int_0^1\{(\sqrt y)^2-y^2\}\dd y
 =\pi\left[\frac{y^2}{2}-\frac{y^3}{3}\right]_0^1
 =\frac{\pi}{6}.
\]
断面を円環として積み重ねるか,円筒の層として重ねるかによって,被積分関数は異なる.しかし,それぞれが立体を重複なく分けているので同じ体積になる.

式\eqref{式:円筒殻の体積}で $x\geq0$ としたのは,異なる $x$ が異なる半径に対応するためである.領域が $y$ 軸をまたぐと,\,$x$ と $-x$ が同じ半径になる.たとえば $-1\leq x\leq1$,\,$0\leq y\leq1$ の長方形を回すと半径 $1$,高さ $1$ の円柱で,体積は $\pi$ である.ところが $2\pi\int_{-1}^1|x|\dd x=2\pi$ は同じ層を二度数えている.絶対値で半径を正にしただけでは,重複は解消しない.



